\documentclass{ximera}


\title{Project 3: Jessica Daniels and Drieyona Lowe}
\author{Jessica Daniels and Drieyona Lowe}

\begin{document}
\begin{abstract}
    Project 3 on kaleidoscope completed by Jessica Daniels and Drieyona Lowe .
\end{abstract}

\maketitle

\section*{Virtual Kaleidoscopes}
\begin{center}
%\vskip 0.125in
\includegraphics[width = 2in]{start.jpg}  \includegraphics[width = 2in]{finish.jpg}
\end{center}

A physical kaleidoscope\footnote{There are many different kinds of kaleidoscopes.  This is just one example.} starts with a pie-shaped image like the one on the left above.  Two mirrors are held perpendicular to the paper with the image.  One mirror is held above the horizontal edge of the image and the other is held above the angled edge of the image.  The two edges and, thus, the two mirrors meet in a $60^{\circ}$ angle.  The horizontal mirror alone would produce the reflection on the left below.  The diagonal mirror alone would produce the reflection on the right below. 

\begin{center}
\includegraphics[width = 2in]{horizontal-reflection.jpg} \includegraphics[width = 2in]{diagonal-reflection.jpg}
\end{center}

\begin{center}
 \includegraphics[width = 4in]{finish.jpg}
  \end{center}
Kaleidoscopic images like the one shown above are formed by reflections, reflections of reflections, reflections of reflections of reflections, and so on.  Each reflection produces a new pie-shaped piece (or sector) of the final image.  Why do you think there are only six pie-shaped pieces or sectors?

We will use two-dimensional vectors $\vec x = \langle{x_1, x_2}\rangle$ to denote the coordinates of each point in the original image.  Mirrors flip or reflect images.  The same point whose original coordinates were denoted by the vector $\vec x = \langle{x_1, x_2}\rangle$ will now have coordinates denoted by a new vector $\vec y = \langle{y_1, y_2}\rangle$ after reflection.  
The new coordinates are determined by a linear transformation

\begin{eqnarray*}
y_1 &=& a_{11} x_1 + a_{12} x_2\\
y_2 &=& a_{21} x_1 + a_{22} x_2\\
\end{eqnarray*}

$$\left[{\begin{array}{c} y_1 \\ y_2 \\ \end{array}}\right] = \left[{\begin{array}{cc} a_{11} & a_{12} \\ a_{21} & a_{22} \\ \end{array}}\right] \left[{\begin{array}{c} x_1 \\ x_2 \\ \end{array}}\right]
\quad {\rm or} \quad \vec y = A \vec x.$$

\begin{enumerate}
    \item Find the linear transformation that reflects an image by a mirror on the $x$-axis.  Express it as a matrix.  Write the matrix below.

$$\left[{\begin{array}{cc}
\\
\ 1 & \ 0 \\
\\
\ 0 & \ -1 \\
\end{array}}\right].$$

\item Find the linear transformation that rotates an image counterclockwise through an angle of $60^{\circ}$.  Express it as a matrix.  Write the matrix below.

$$\left[{\begin{array}{cc}
\\
\ 1/2 & -\sqrt{3}/2 \\
\\
\sqrt{3}/2& \ 1/2 \\
\end{array}}\right].$$


\item Find the linear transformation that reflects an image by a mirror on the diagonal line that makes an angle of $60^{\circ}$ with the $x$-axis.  Express it as a matrix.  Write the matrix below.

$$\left[{\begin{array}{cc}
\\
\ -1/2 & \sqrt{3}/2 \\
\\
\sqrt{3}/2 & \ 1/2 \\
\end{array}}\right].$$

\vspace{8pt}


\noindent Explain how you developed this matrix.


We started by figuring out how to represent the reflection of an image across a diagonal line. We found that this can be represented by a  linear transformation that reflects vertically creating the point $(x,y)$; when reflected horizontally or across the $y$-axis, it results in the coordinate $(x,y)$. To find the matrix, we looked for a $2\times 2$ matrix
$A=\left[
\begin{array}{cc}
   a & b \\
   c  & d
\end{array}
\right].$ when multiplied by $ \left[ \begin{array}{cc}
     x\\
     y 
\end{array} \right] $, produces the vector  $ \left[ \begin{array}{cc}
     x\\
     -y 
\end{array} \right] $. 
Therefore, for this to be true for any $x$ and $y$,

$ax+by=x$, implies $a=1$ and $b=0$, as well as $cx+dy=-y$, which implies $c = 0$ and $d= -1$.  

For question 2, it asks us to rotate a point counterclockwise by an angle $\theta$ around a 2D plane, where we applied the standard rotation matrix. 

$$R(\theta)= \left[ \begin{array}{cc}
   \cos\theta  & -\sin\theta \\
   \sin\theta  & \cos\theta
\end{array} \right]$$
where $\theta= 60^\circ$.

For question 3, we did similar process to question 2, but considering the positive x-axis. This transformation matrix = $$M(\alpha)= \left[ \begin{array}{cc}
   \cos(2\alpha)  & \sin(2\alpha) \\
    \sin(2\alpha) & -\cos(2\alpha)
\end{array}{} \right]$$

where $\alpha= 60^\circ$  and $2\alpha= 120^\circ$.

\end{enumerate}
\end{document}  